% Аннотация (рус.)
\specialsection{Аннотация}

Аннотация

Выпускная квалификационная работа посвящена адаптивному отказоустойчивому хранению данных в распределённой системе с дозаписью. Цель работы состоит в разработке модели и архитектурного подхода, при котором схема хранения может изменяться в фоне в течение жизненного цикла данных, а не выбираться один раз при записи объекта. В рамках работы проанализированы подходы на основе репликации, кодов стираний и температурных политик, предложена модель гибридного хранения с репликами, локальными и глобальными проверочными блоками, описаны допустимые переходы между схемами и модель оценки их полезности с учётом активности данных, занятого места, ожидаемой стоимости чтения и стоимости фонового перекодирования. Также разработана экспериментальная система с фоновым перекодировщиком и планировщиками. Эксперименты в сценариях охлаждения данных, неравномерной нагрузки чтения и отказов узлов показывают, что адаптивное перекодирование позволяет эффективнее использовать дисковое пространство по сравнению с репликацией и сохранять более дешёвое чтение для более активных или рискованных участков по сравнению с постоянным использованием кодов стираний.

Ключевые слова: распределённое хранение; коды стираний; репликация; температура данных; фоновое перекодирование.

\pagebreak
